\chapter*{概要}

\addcontentsline{toc}{chapter}{概要}

都市の状況把握において，センシングデータや統計データは客観性が高い一方，更新頻度の低さや設置コストの高さ，市民の主観的・情動的反応の把握困難といった課題がある．
これに対し，SNS上のコメントは市民の反応を捉える「ヒューマンセンサ」として注目されてきたが，主要プラットフォームであるX（旧Twitter）ではAPI制限の厳格化により，継続的なデータ収集が困難となっている．
また，既存研究の多くは災害発生後の限定的な期間を対象とした分析にとどまり，平常時から災害時まで一貫して扱える持続的な分析基盤は十分に整備されていない．

本研究では，新たなSNSコメント収集ソースとしてYouTubeライブ配信コメントに着目し，災害時における市民の反応を捉える分析基盤を構築した．
特に，日本全国の地震情報をリアルタイムで配信するYouTubeライブチャンネルに集積するコメントを対象とし，平常時から災害時まで継続的にデータを収集・蓄積・分析可能なシステムを提案した．

提案システムは，データ収集層，蓄積層，分析層，可視化層の4層から構成される．
データ収集層では，YouTube Data API v3を用いたポーリング方式により，AWS Lambdaのサーバレス環境で低コストな収集を実現した．
データ蓄積層では，BigQueryを採用し時間パーティショニングによる効率的なクエリ処理を可能とし，Infrastructure as Code（Terraform）により環境の再現性とスケーラビリティを確保した．
分析層では，BERT及びLUKE-WRIMEによる感情分析を独立したモジュールとして実装し，拡張性を持たせた．
可視化層では，Looker Studioによる探索的データ分析とNext.js製Webダッシュボードによる即時的な状況把握を両立させた．

約9ヶ月間（2025年2月〜10月）の運用により，約48万件のデータを収集し，システムの安定稼働を確認した．
収集データの分析から，YouTubeライブコメントは平均約15文字と短文による即応的な反応が特徴であり，地震発生後にコメント数が急増する明確な反応パターンが確認された．
また，コメント投稿には時間帯や曜日による傾向があり，特定のユーザによる集中的な投稿といった特性も明らかとなった．
感情分析の結果，地震規模とコメント数，感情の関係性が示され，災害時の市民感情の変化を定量的に捉えることが可能であることを確認した．

本研究により，YouTubeライブ配信コメントは，従来のSNSデータとは異なる特性を持つ新たなヒューマンセンサデータソースとして，災害情報収集において一定の有効性を持つことが示された．
提案した分析基盤は，低コストかつスケーラブルな構成により継続的な運用が可能であり，SNSを活用した都市分析基盤の高度化に資する知見を提供するものである．
